\documentclass[12pt,aspectratio=43]{beamer}

% 自訂字體的封包
\usepackage{fontspec}

% 導入圖形與表格的封包
\usepackage{graphicx}
\graphicspath{{./20251029_Meeting}}

\usepackage{booktabs}

% 排列多個子圖形的封包
\usepackage{subfigure} 

 % 提供進階數學環境與公式排版
\usepackage{mathtools, amsmath, amsfonts, amsthm, latexsym} 

%% 設定中文字體
\usepackage{xeCJK}
\setCJKmainfont[AutoFakeBold=3.5 , AutoFakeSlant=0.2]{標楷體}
\setmainfont{Times New Roman}
\XeTeXlinebreaklocale "zh"
\XeTeXlinebreakskip = 0pt plus 1pt

% 自訂字體顏色的封包
\usepackage{color} 

% 圖表自動編號的封包
\usepackage{caption} 

% 允許表格的一格能多列呈現的封包
\usepackage{multirow} 

% 可指定表格排版的封包
\usepackage{array}

% 翻轉表格的封包
\usepackage{lscape} 

% 序列標號
\usepackage{enumerate} 

%% 設定自動編號
\setbeamertemplate{caption}[numbered]

%% 自訂顏色
\definecolor{Myblue}{RGB}{57,108,144}
\definecolor{Default_Blue}{RGB}{52,51,171}

% (動畫)字會若隱若現的顯示
\beamersetuncovermixins{\opaqueness<1->{45}}{\opaqueness<1->{95}} 

% 設定中文的標籤
\renewcommand{\figurename}{圖} 
\renewcommand{\tablename}{表} 

% 排版形式 
\mode<presentation>{
\usetheme{Berlin}
% enumerate 及 item 的形狀
\setbeamertemplate{itemize items}[triangle]
% 調整外框形式的字體大小
\setbeamerfont{headline}{size=\scriptsize}
\setbeamerfont{footline}{size=\scriptsize}
% 取消右下方的跳轉工具列
\setbeamertemplate{navigation symbols}{} 
% 自定義2：清除外框下緣但僅出頁碼
\setbeamertemplate{footline}[page number] 
% 顏色主題
\usecolortheme{dove}
% 設定頁面邊界
\setbeamersize{text margin left=0.6cm, text margin right=0.6cm}
\special{papersize=\the\paperwidth,\the\paperheight}
\providecommand{\tabularnewline}{\\}
}

\title{Meeting 1029}
\author{黃丞暘}
\institute{國立東華大學應數統計所}
\date{October 29, 2025}

\begin{document}

% 封面
\begin{frame}
  \titlepage
\end{frame}

% 目錄
\begin{frame}{\textbf{目錄}}
  \tableofcontents
\end{frame}

% 研究動機 
\section{研究動機}

\begin{frame}{\textbf{研究動機}}
\linespread{1.5}
  Semiconductor Manufacturing Process
  \begin{itemize}
    \item Robust performance can be trust
  \end{itemize}
  \begin{itemize}
    \item High dimension
  \end{itemize}
  \begin{itemize}
    \item Find the $y = \hat{f}(x)$
  \end{itemize}
\end{frame}

% 研究方法
\section{研究方法}
\begin{frame}{\textbf{研究方法}}
\linespread{1.5}
  \begin{itemize}[]
    \item 前處理
      \begin{enumerate}[]
        \item vote in features selection
      \end{enumerate}
    \item 機器學習模型
      \begin{enumerate}[]
        \item XGboost, RandomForest, DecisionTree, SVM
      \end{enumerate}
    \item 
  \end{itemize}
\end{frame}

% --- 參考資料 ---
\section{參考資料}

\begin{frame}{\textbf{參考資料}}
\linespread{1.5}

\footnotesize

  \begin{thebibliography}{99} 
    \bibitem[Li-Fu Lai (2022)]{中文文獻範例} % 引注時顯示的名稱
      Li-Fu Lai (2022) % 作者 (年份)
      \newblock 使用KNN和決策樹方法對填補缺失值的製造數據進行失效預測和原因追蹤分析 % 文章名稱
  \end{thebibliography}
\end{frame}

% --- 結束 ---
\section*{結束}

\begin{frame}[plain] % 產生無外框的頁面
	\begin{center}
		{\Huge The End}
		\bigskip\bigskip 

		{\LARGE Questions  \&  Comments}
	\end{center}
\end{frame}
\end{document}ㄒ
